\begin{problem}[A concrete contravariant map: squaring]
Let $k$ be algebraically closed and consider
\[
\varphi:k[t]\to k[x],\qquad t\mapsto x^2.
\]
Describe the induced map
\[
\varphi^\ast:\Spec k[x]\to\Spec k[t]
\]
on all prime ideals.  Discuss its behavior on closed points in characteristic different from
$2$, and in characteristic $2$.
\end{problem}

\paragraph{Why this problem matters.}
This is a fully explicit example where a polynomial map is encoded by contraction of primes.

\paragraph{Useful ideas before starting.}
The primes of each polynomial ring are $(0)$ and linear maximal ideals.

\paragraph{Guided hints.}
Compute the contraction of $(x-a)$ by evaluating a polynomial $f(t)$ at $t=a^2$.

\begin{solution}
The zero prime contracts to zero because $\varphi$ is injective:
\[
\varphi^{-1}((0))=(0).
\]
For a closed point $(x-a)$,
\[
f(t)\in\varphi^{-1}(x-a)
\iff
f(x^2)\in(x-a)
\iff
f(a^2)=0.
\]
Hence
\[
\varphi^\ast((x-a))=(t-a^2).
\]

If $\operatorname{char}k\neq2$, a nonzero closed point $(t-b)$ has two preimages whenever
$b\neq0$, corresponding to the two square roots $\pm\sqrt b$, while $b=0$ has the single
preimage $(x)$.  Since $k$ is algebraically closed, every $b$ has a square root.

If $\operatorname{char}k=2$, the map $a\mapsto a^2$ is injective on a field and surjective
because $k$ is algebraically closed, so every closed point has a unique closed-point preimage.
The generic point maps to the generic point in all characteristics.
\end{solution}

\paragraph{What happened mathematically.}
The same ring homomorphism can have visibly different point-fiber behavior depending on the
characteristic, even though the contraction formula is uniform.

\paragraph{Extensions.}
Replace $x^2$ by $x^n$ and analyze the effect of the characteristic dividing $n$.

\paragraph{Provenance.}
New canonical example/problem illustrating the contravariant rule.
